\textbf{COLUNI 2019 -- Questão 28}\par\medskip O vírus da AIDS ou HIV é responsável pela síndrome da imunodeficiência adquirida. Atualmente, o tratamento de pacientes infectados pelo HIV se dá por um coquetel de medicamentos distribuídos gratuitamente pelo Sistema Único de Saúde (SUS). Esse tratamento permite levar uma vida praticamente normal por inibir a replicação do vírus, não sendo, porém, capaz de eliminá-lo totalmente do organismo. Pacientes que abandonam o tratamento poderão desenvolver os sintomas da doença e vir a óbito.

\par

Assinale a alternativa que apresenta CORRETAMENTE a explicação sobre o que acontece com pacientes que são infectados pelo HIV, mas não tomam o coquetel:\begin{enumerate}
\item O sistema reprodutor deixa de produzir gametas em função da baixa produção de hormônios sexuais.
\item O sistema imunológico deixa de combater patógenos em função da destruição das células de defesa.
\item O sistema circulatório deixa de oxigenar os órgãos em função da redução da atividade cardíaca.
\item O sistema digestivo deixa de digerir os alimentos em função da perda dos movimentos peristálticos.\end{enumerate}
